\section{Numerical rank and tangential limits}\label{sec:limits}
We now compare a boundary restriction of the extended highest line
with the limiting projective line map. The Hodge-theoretic input
identifies the limiting graded line. The analytic rank principle
then decides what would happen if its tangential rank did not drop.
For a nef line $J$ on a smooth projective $t$-fold we use
\[
 \nu(J)=\max\{a\in\{0,\ldots,t\}:c_1(J)^aH^{t-a}>0\},
\]
where $H$ is ample. A line on a point has numerical dimension zero.

Let $U$ and its variation be as in Section~\ref{sec:rational}.
Take a smooth projective compactification $T$ of $U$ with reduced
simple normal crossings boundary $D$, and suppose that all local
monodromies along $D$ are unipotent. Denote the Deligne extension by
$\overline V$ and the Schmid extension of the Hodge line by
$J=F^m\overline V$. We use the following standard inputs.

\begin{theorem}[Extension and positivity inputs]\label{altord:h:inputs}
Under these hypotheses:
\begin{enumerate}
\item The extended Hodge filtration is by subbundles of $\overline V$.
For a morphism $g:(T',D')\to(T,D)$ of smooth projective varieties
with reduced simple normal crossings boundaries and
$g(T'\setminus D')\subset U$, the extended Hodge line of $g^*V$
is $g^*J$. These extensions commute with direct sums and tensor products.
If the monodromy logarithms of some boundary components vanish, the
variation extends purely across those components, including their
crossings away from the remaining boundary.
\item The line $J$ is nef. It is big precisely when the top
Kodaira--Spencer map
\[
 T_U\longrightarrow
 \Hom(F^mV_{\mathcal O},F^{m-1}V_{\mathcal O}/F^mV_{\mathcal O})
\]
is generically injective, equivalently when $\ell$ is generically
immersive.
\end{enumerate}
The first assertions are the unipotent extension and nilpotent-orbit
theorems of \cite{Schmid,CKS}; functoriality is also recorded in
\cite[Lemmas~2.16 and~5.6]{BFMT}. The second assertion is
\cite[Lemma~6.17]{BakkerBrunebarbeTsimerman}.
\end{theorem}

The equivalence in the second assertion uses Griffiths transversality:
the differential of the line map is the indicated Kodaira--Spencer
map followed by the inclusion of $F^{m-1}/F^m$ in $V_{\mathcal O}/F^m$.

\begin{proposition}[Numerical rank]\label{altord:h:numerical-rank}
If $r=\rk_{\mathrm{gen}}d\ell$, then
\begin{equation}\label{altord:h:eq-numerical-rank}
 \nu(J)=r.
\end{equation}
\end{proposition}

\begin{proof}
Put $t=\dim T$, and fix an ample divisor $H$ on $T$.
For $1\le a\le t$, choose a sufficiently general smooth complete
intersection $C_a\subset T$ of dimension $a$, using a sufficiently
large multiple of $H$. It meets the boundary transversely. At general
points of $C_a\cap U$, the restriction of $d\ell$ to $T_{C_a}$ has
rank $\min(a,r)$. To see the required genericity, at a point where
$d\ell$ has rank $r$, the $a$-planes with this property form a nonempty
open subset of the tangent Grassmannian. A sufficiently ample linear
system prescribes the required first-order tangent conditions; Bertini
then gives the assertion for general complete intersections.
By Theorem~\ref{altord:h:inputs}, $J|_{C_a}$ is big exactly when $a\le r$.
Since it is nef, this is equivalent to
$(J|_{C_a})^a>0$, and hence to $c_1(J)^aH^{t-a}>0$.
This proves the formula, also when $r=0$; dimension zero is immediate.
\end{proof}

\subsection{The limiting line along a divisor}

We recall the boundary input in the form needed here. Let $E$ be a
component of $D$, let
$E^\circ=E\setminus\Supp(D-E)$, and let $N$ be its monodromy logarithm.
The monodromy weight filtration $W(N)$ is centred at $w$.
The limiting filtration induces polarized pure variations on its graded
pieces over $E^\circ$, with unipotent monodromy at the remaining boundary.
If $k$ is the unique weight carrying the deepest line, then $J|_E$ is
the Schmid extension of the top line of $\Gr^{W(N)}_k$.
These statements include compatibility at deeper strata;
see \cite[\S2.5.2--\S2.5.3 and Lemma~2.16(3)]{BFMT}.
For a single $N$, the relevant primitive summand is polarized by the
form induced by $Q(-,N^{k-w}-)$ when $k\ge w$.

The local expression below is Schmid's nilpotent-orbit normal form
\cite[Theorem~4.12]{Schmid}; its compatibility at crossings uses
\cite{CKS}. The rank comparison is the consequence needed here.

\begin{proposition}[Boundary line]\label{altord:h:boundary-line}
On a transverse coordinate chart $(u,v)$ about a general point of $E$,
with $E=(u=0)$, write a generator of the lifted Hodge line as
\begin{equation}\label{altord:h:nilpotent-expression}
 e^{zN}a(u,v),\qquad z=\frac{\log u}{2\pi i},
\end{equation}
where $a$ is holomorphic across $u=0$ and $a(0,v)\ne0$.
After shrinking the chart there is an integer $\ell\ge0$ such that
\begin{equation}\label{altord:h:primitive-line}
 N^{\ell+1}a(0,v)=0,\qquad N^\ell a(0,v)\ne0,
\end{equation}
and the line belongs to weight $w+\ell$ in the limiting mixed Hodge
structure. Moreover,
\begin{equation}\label{altord:h:boundary-rank}
 \nu(J|_E)\le
 \rk_{\mathrm{gen}}d\bigl(v\longmapsto[N^\ell a(0,v)]\bigr).
\end{equation}
The map on the right is formed in the chosen flat coordinates.
If $N=0$, take $\ell=0$; the variation then extends purely near the
general point of $E$.
\end{proposition}

\begin{proof}
Put $a_0=a(0,v)$. In the Deligne bigrading of the limiting mixed
Hodge structure, $F^m$ is a single one-dimensional summand $I^{m,q}$:
all first indices are at most $m$, and $\dim F^m=1$.
Set $\ell=m+q-w$. The primitive decomposition on weight $w+\ell$
expresses each summand as a sum of terms
\[
 N^j P_{w+\ell+2j},\qquad j\ge\max(0,-\ell),
\]
where $P_{w+b}$ denotes the primitive part for $b\ge0$.
A contribution to Hodge type $(m,q)$ with $j>0$ would come from
type $(m+j,q+j)$, impossible because its first index exceeds $m$.
Thus $j=0$, $\ell\ge0$, and the top line is primitive.
Hard Lefschetz for the monodromy weight filtration gives nonvanishing
of its image under $N^\ell$ and vanishing under $N^{\ell+1}$ on the
associated graded. Since $N$ has bidegree $(-1,-1)$ in the Deligne
bigrading, these statements hold for $a_0$ itself: a vector in a single
Deligne summand maps injectively to its weight-graded piece.
The integer $\ell$ is constant on a sufficiently small open stratum,
so the identities hold throughout the chosen $v$-chart.

Let $J_E^{\mathrm{gr}}$ be the extended top line of the pure variation
on $\Gr^{W(N)}_{w+\ell}$ along $E^\circ$.
The boundary extension input identifies it with $J|_E$.
Projecting $N^\ell a_0$ to $\Gr^{W(N)}_{w-\ell}$ yields the image
of this graded top line under the flat isomorphism
\[
 N^\ell:\Gr^{W(N)}_{w+\ell}\longrightarrow
          \Gr^{W(N)}_{w-\ell}.
\]
Consequently the projective map of the graded line has rank no greater
than $v\mapsto[N^\ell a_0(v)]$: projection is a fixed linear map on
the flat chart, defined near the line in question.
Apply Proposition~\ref{altord:h:numerical-rank} to the graded variation on
$E^\circ$ and its smooth projective compactification $E$.
This proves \eqref{altord:h:boundary-rank}.
\end{proof}

\subsection{A local rank principle}

The boundary-line proposition has reduced the numerical question to a
projective differential. We next prove the local assertion that will
force a symmetry to fix the interior whenever the limiting rank is
as large as the interior rank. This analytic fact requires control of tangential derivatives
only. A sector below has the form
$S_\epsilon=\{0<|u|<\epsilon,\ \alpha<\arg u<\beta\}$ with a fixed
branch of the argument and finite opening $\beta-\alpha<2\pi$.

\begin{lemma}[Rank and a tangential limit]\label{altord:h:rank-principle}
Let $B\subset\C^n$ be a polydisc and let
$f:S_\epsilon\times B\to M$ be a holomorphic map to a complex
manifold. Suppose $f(u,\cdot)$ converges uniformly on compact subsets
of $B$, as $u\to0$ throughout the sector, to a holomorphic map $f_0$.
Assume that the full generic differential rank of $f$ is at most $r$
and that $\rk df_0(v_0)=r$ at some $v_0\in B$.
Then, after shrinking the sector and taking neighborhoods
$v_0\in B'\Subset B_0\subset B$,
\[
 f(S_{\epsilon'}\times B')\subset f_0(B_0),
\]
where the indicated local image of $f_0$ is an embedded complex
submanifold of dimension $r$. In particular, every fixed closed analytic
subset containing this local image also contains the nearby image of $f$.
\end{lemma}

\begin{proof}
Use a common coordinate chart in $M$, shrinking $B$ and the sector.
Cauchy estimates give convergence of all derivatives in the $B$
variables on smaller polydiscs. The $(r+1)$-minors of $df$ vanish
identically, so those of $df_0$ vanish as well. Thus $f_0$ has constant
rank $r$ near $v_0$. The holomorphic constant rank theorem supplies
source coordinates $v=(x,y)\in\C^r\times\C^{n-r}$ and target
coordinates in which $f_0(x,y)=(x,0)$.

Write $f=(F,G)$ in these coordinates, with
$F(u,x,y)=x+E(u,x,y)$. On fixed smaller polydiscs, both $E$ and
$\partial_xE$ tend uniformly to zero. Choose nested $x$-balls, a
smaller ball $W$ for $w$, and a ball $Y$ for $y$. For sufficiently
small $\epsilon'$, the maps
\[
 x\longmapsto w-E(u,x,y)
\]
send a fixed closed $x$-ball into itself and have Lipschitz constant
less than $1/2$, uniformly for
$(u,w,y)\in S_{\epsilon'}\times W\times Y$.
The contraction theorem gives a unique solution
$x=\chi(u,w,y)$ of $F(u,x,y)=w$ on this common product.
The holomorphic inverse function theorem and uniqueness make $\chi$
holomorphic in all interior parameters; moreover $\chi\to w$
uniformly on smaller products.

In the new source coordinates the map is
\[
 f(u,\chi(u,w,y),y)=(w,H(u,w,y)).
\]
The $w$-derivatives already have rank $r$.
The rank bound therefore forces $\partial_uH=0$ and
$\partial_yH=0$: either derivative, if nonzero, would give an
additional image vector with zero first $r$ coordinates.
Since the sector and $Y$ are connected, $H$ depends only on $w$.
Its limit as $u\to0$ is zero, hence $H=0$.
Shrinking the original source neighborhood so that its $F$-values
belong to $W$ proves the assertion. If $r=0$, the rank bound directly
makes $f$ constant on the connected source, and the same conclusion
follows.
\end{proof}

\begin{lemma}[Nilpotent tangential convergence]\label{altord:h:nilpotent-limit}
Let $N$ be a nilpotent endomorphism of a finite-dimensional complex
vector space, and let $a(u,v)$ be holomorphic on $\Delta\times B$.
Suppose, for a fixed $\ell\ge0$, that
$N^{\ell+1}a(0,v)=0$ and $N^\ell a(0,v)\ne0$ throughout $B$.
On a sector with the above choice of $z=\log u/(2\pi i)$,
\begin{equation}\label{altord:h:eq-nilpotent-limit}
 [e^{zN}a(u,v)]\longrightarrow[N^\ell a(0,v)]
\end{equation}
normally on sufficiently small $v$-charts, with all tangential
derivatives. Thus Lemma~\ref{altord:h:rank-principle} applies whenever the
rank of this limit equals an upper bound for the full rank of the
line map. Its nearby image then lies in $\mathbb P(\ker N)$.
\end{lemma}

\begin{proof}
Choose $k$ with $N^{k+1}=0$, and write
$a(u,v)=a_0(v)+u q(u,v)$. Set $b(v)=N^\ell a_0(v)/\ell!$.
Direct expansion gives
\begin{align*}
 z^{-\ell}e^{zN}a(u,v)
 &=b(v)+\sum_{j<\ell}\frac{z^{j-\ell}}{j!}N^ja_0(v)
       +u\sum_{j=0}^k\frac{z^{j-\ell}}{j!}N^jq(u,v).
\end{align*}
For every fixed tangential derivative, uniformly on compact subsets
of $B$, the remainder is
\[
 O(|z|^{-1})+O(|u|\,|z|^{k-\ell});
\]
the first term is absent when $\ell=0$.
On the sector $|z|$ is comparable to $|\log|u||$, so both terms tend
to zero. Choose a linear functional nonzero on $b(v_0)$ and shrink
the $v$-chart. Dividing by that functional gives convergence with
all tangential derivatives in a fixed affine projective chart,
proving \eqref{altord:h:eq-nilpotent-limit}. Finally $Nb=0$, so the limiting
image belongs to the fixed projective linear subspace
$\mathbb P(\ker N)$, and the last assertion follows from
Lemma~\ref{altord:h:rank-principle}.
\end{proof}

\begin{remark}\label{altord:h:scaling}
All numerical-rank statements are unchanged on replacing a Hodge line
by a positive rational multiple of its class in $\Pic(T)\otimes\Q$.
The local statements also apply with $N=0$, when the line map extends
holomorphically across the divisor.
\end{remark}
